% ============================================================================
%   CHAPTER 5 — RESULTS AND ANALYSIS
%   Aivancity: 10-15 pages — MUST BE COMPACT
%   Strategy: main results in body, all details in appendices
%
%   REORGANIZED to integrate mode collapse and non-trivial case analysis
% ============================================================================

\chapter{Results and Analysis}
\label{chap:results}

% ----------------------------------------------------------------------------
\section{Chapter Overview}
\label{sec:res-overview}

% [TODO — 0.5 page]
% Reader's guide:
% - Section 5.2: pre-LLM audit and expert-GPT4 agreement
% - Section 5.3: main calibration audit (H1)
% - Section 5.4: RLHF effect (H2), main pre-registered contribution
% - Section 5.5: the enriched-prompt paradox
% - Section 5.6: mode collapse (novel finding, 5 of 12 configurations)
% - Section 5.7: performance on truly discriminative cases
%     (novel finding, aggregate metrics hide complete failure)
% - Section 5.8: inclusion/exclusion asymmetry
% - Section 5.9: model-specific evasion biases
% - Section 5.10: ordinal validity (H3) and selective classification (H4)
% - Section 5.11: post-hoc recalibration and its cross-validation stability
% - Section 5.12: conformal prediction and inter-model consensus
% - Section 5.13: baselines and composite utility score
% - Section 5.14: summary of results and findings

% ----------------------------------------------------------------------------
\section{Pre-LLM Audit and Expert Agreement}
\label{sec:res-audit}

% [TODO — 1 page]
% Kappa 0.816 expert vs GPT-4
% Ordinal agreement metrics
% First original finding: GPT-4 "not applicable" evasion bias (ratio 2.16)

\begin{finding}[F1: GPT-4's Selective Evasion Bias]
    GPT-4 over-uses the ``not applicable'' label (ratio 2.16 vs experts) but
    uses ``not enough information'' with a ratio of 0.99. This asymmetric
    pattern reveals a selective evasion, unreported in the original
    TrialGPT paper.
\end{finding}

% ----------------------------------------------------------------------------
\section{Calibration Audit: H1 Supported Across All Models}
\label{sec:res-h1}

% [TODO — 2 pages]
% The 12 ECE values with bootstrap CIs
% Main summary table
% Reliability diagram (the flagship figure)

\begin{table}[htbp]
    \centering
    \caption{Calibration metrics across the 12 configurations.}
    \label{tab:res-calibration}
    \begin{tabular}{@{}lccc@{}}
        \toprule
        \textbf{Configuration} & \textbf{ECE} & \textbf{95\% CI} & \textbf{Brier} \\
        \midrule
        Mistral-Base A     & 0.133 & [0.104, 0.163] & 0.258 \\
        Mistral-Instruct A & 0.322 & [0.297, 0.350] & 0.311 \\
        Mistral-Base B     & 0.189 & [0.161, 0.218] & 0.269 \\
        Mistral-Instruct B & 0.253 & [0.229, 0.287] & 0.353 \\
        Llama-Base A       & 0.124 & [0.100, 0.148] & 0.187 \\
        Llama-Instruct A   & 0.510 & [0.483, 0.535] & 0.456 \\
        Llama-Base B       & 0.151 & [0.123, 0.177] & 0.225 \\
        Llama-Instruct B   & 0.602 & [0.575, 0.628] & 0.558 \\
        Meditron A         & 0.182 & [0.153, 0.212] & 0.263 \\
        Meditron B         & 0.231 & [0.204, 0.259] & 0.283 \\
        BioMistral A       & 0.155 & [0.126, 0.182] & 0.234 \\
        BioMistral B       & 0.194 & [0.166, 0.222] & 0.244 \\
        \bottomrule
    \end{tabular}
\end{table}

\begin{figure}[htbp]
    \centering
    % \includegraphics[width=0.9\textwidth]{reliability_flagship.pdf}
    \caption{Reliability diagrams of all 12 configurations (Prompt A left,
        Prompt B right). All curves fall well below the perfect calibration
        diagonal.}
    \label{fig:res-reliability}
\end{figure}

% ----------------------------------------------------------------------------
\section{The Effect of RLHF: H2 Supported (Main Contribution)}
\label{sec:res-h2}

% [TODO — 2 pages]
% The 4 permutation tests: Mistral A, Mistral B, Llama A, Llama B
% All 4 significant after Bonferroni
% Emphasis: our original contribution, extending Kadavath 2022 to clinical

\begin{finding}[F2: RLHF Systematically Degrades Calibration]
    Across two independent model families (Mistral-7B and Llama-3.1-8B) and
    two independent prompt designs, RLHF-aligned Instruct versions exhibit
    significantly worse calibration than their Base counterparts
    (all $p < 0.003$ with Bonferroni correction, $\Delta\ece$ ranging from
    +0.064 to +0.451).
\end{finding}

% Interpretation: RLHF transforms modesty into overconfidence.
% Show table comparing mean overconfidence:
%   Mistral-Base: −0.189 → Mistral-Instruct: +0.248
%   Llama-Base: +0.128 → Llama-Instruct: +0.598

% ----------------------------------------------------------------------------
\section{The Enriched-Prompt Paradox}
\label{sec:res-prompt}

% [TODO — 1.5 pages]
% Only Mistral-Instruct benefits from Prompt B
% 5 out of 6 models are DEGRADED by the expert prompt

\begin{finding}[F3: The Enriched-Prompt Paradox]
    Contrary to the intuition that expert prompt design improves calibration,
    prompts enriched with definitions, examples, and exclusion rules
    \emph{degrade} calibration in 5 of 6 models. Only the strongly
    RLHF-aligned Mistral-Instruct benefits.
\end{finding}

% ----------------------------------------------------------------------------
\section{Mode Collapse in Medical and RLHF-Aligned LLMs}
\label{sec:res-mode-collapse}

% [TODO — 2 pages — NEW SECTION, one of the flagship findings]
%
% Content:
% 1. Motivation: standard calibration metrics assume the model produces
%    diverse predictions. What if a model always predicts the same label?
% 2. Distribution analysis: 5 out of 12 configurations predict "not enough
%    information" in more than 97% of cases:
%    - BioMistral-B: 100.0%
%    - Llama-Instruct-B: 99.7%
%    - BioMistral-A: 98.2%
%    - Llama-Instruct-A: 97.8%
%    - Llama-Base-B: 97.2%
% 3. Consequence: their "reasonable" ECE (0.15–0.60) reflects proximity to
%    a naïve NEI baseline (ECE 0.21), not genuine predictive quality.
% 4. Meditron partial collapse: predicts only "not excluded" (65%) and NEI (35%).
% 5. Simulation experiment: forcing models to choose between two decisive
%    labels does NOT restore reasoning — collapse simply shifts to the
%    majority label per criterion type.

\begin{finding}[F7: Mode Collapse in 5 of 12 Configurations]
    Five open-source configurations --- both BioMistral variants, all
    Llama-Instruct variants, and Llama-Base with prompt B --- predict
    ``not enough information'' in more than 97\% of cases. Their apparently
    reasonable ECE (0.15–0.60) reflects proximity to a trivial NEI baseline
    rather than genuine probabilistic quality.
\end{finding}

\begin{finding}[F8: Mode Collapse Is Structural, Not Superficial]
    When forced to choose between the two decisive labels of each criterion
    type by removing evasion candidates from the softmax, collapsed models
    do not recover reasoning ability. Instead, collapse migrates from
    ``not enough information'' to the majority label per criterion type
    (``not excluded'' for exclusion criteria, or ``excluded'' for one
    Llama-Base variant). Evasion behaviors mask absence of clinical
    reasoning rather than strategic uncertainty communication.
\end{finding}

% ----------------------------------------------------------------------------
\section{Performance on Truly Discriminative Cases}
\label{sec:res-non-trivial}

% [TODO — 2 pages — NEW SECTION, second flagship finding]
%
% Content:
% 1. Motivation: 76.6% of cases have ground truth in {not excluded, NEI};
%    the remaining 23.4% require true clinical discrimination.
% 2. On these 238 truly discriminative cases:
%    - Meditron A: 0.0% accuracy (predicts only NE/NEI)
%    - Meditron B: 0.0% accuracy
%    - BioMistral A: 0.4% accuracy
%    - BioMistral B: 0.0% accuracy
%    - Llama-Base B: 1.3% accuracy
%    - Llama-Instruct B: 1.3% accuracy
%    - Mistral-Instruct A: 71.9% accuracy (only truly usable)
%    - Mistral-Base A: 61.3%
% 3. Complete reversal of the ranking: Meditron, apparent leader in global
%    accuracy (64.8%), is actually WORSE on the cases that matter.
% 4. Only Mistral-Instruct A reaches clinically usable performance
%    (accuracy 71.9%, ECE 0.054 on non-trivial subset).

\begin{finding}[F9: Global Accuracy Hides Complete Failure on Discriminative Cases]
    On the 238 truly discriminative cases (23.4\% of the dataset), Meditron
    achieves 0\% accuracy despite its 64.8\% global accuracy; BioMistral
    achieves 0.0\% and 0.4\% accuracy for its two prompts. The apparent
    superiority of medical LLMs disappears on the cases requiring genuine
    clinical decisions. Only Mistral-Instruct with prompt A reaches
    acceptable clinical performance (71.9\% accuracy, ECE 0.054).
\end{finding}

\begin{table}[htbp]
    \centering
    \caption{Global vs.\ non-trivial accuracy across configurations. The
        performance drop reveals which models rely on the majority class
        distribution.}
    \label{tab:res-non-trivial}
    \small
    \begin{tabular}{@{}lccc@{}}
        \toprule
        \textbf{Configuration} & \textbf{Global acc.} & \textbf{Non-trivial acc.} & \textbf{Drop} \\
        \midrule
        Meditron A         & 0.648 & 0.000 & $-$0.847 \\
        Meditron B         & 0.630 & 0.000 & $-$0.824 \\
        BioMistral B       & 0.289 & 0.000 & $-$0.377 \\
        BioMistral A       & 0.296 & 0.004 & $-$0.381 \\
        Llama-Instruct B   & 0.292 & 0.013 & $-$0.365 \\
        Llama-Base B       & 0.308 & 0.013 & $-$0.386 \\
        Llama-Instruct A   & 0.290 & 0.021 & $-$0.351 \\
        Llama-Base A       & 0.230 & 0.265 & $+$0.046 \\
        Mistral-Base B     & 0.586 & 0.416 & $-$0.222 \\
        Mistral-Instruct B & 0.480 & 0.500 & $+$0.026 \\
        Mistral-Base A     & 0.500 & 0.613 & $+$0.149 \\
        Mistral-Instruct A & 0.308 & \textbf{0.719} & $+$0.536 \\
        \bottomrule
    \end{tabular}
\end{table}

% ----------------------------------------------------------------------------
\section{Systematic Inclusion vs Exclusion Asymmetry}
\label{sec:res-asymmetry}

% [TODO — 1 page]
% Stratified ECE by criterion_type
% Almost all models: much worse on exclusion than inclusion
% Meditron reverses the pattern

\begin{finding}[F4: Inclusion vs Exclusion Asymmetry]
    LLM calibration on exclusion criteria is 2--3 times worse than on
    inclusion criteria across all generalist models. \model{Meditron} is
    the sole exception, showing the reverse pattern --- likely due to its
    clinical guidelines training data.
\end{finding}

% ----------------------------------------------------------------------------
\section{Model-Specific Evasion Biases}
\label{sec:res-evasion}

% [TODO — 1 page]
% GPT-4: over-uses "not applicable" (ratio 2.16)
% Open-source LLMs: over-use "not enough information" (ratio up to 3.46)

\begin{finding}[F5: Divergent Evasion Strategies Across LLMs]
    Whereas GPT-4 evades via the ``not applicable'' label (ratio 2.16), open-
    source LLMs consistently evade via ``not enough information'' (ratio 3.3
    to 3.5 for BioMistral and Llama-Instruct). Evasion is systematic but
    channel-specific.
\end{finding}

% ----------------------------------------------------------------------------
\section{Ordinal Validity (H3) and Selective Classification (H4)}
\label{sec:res-h3h4}

\subsection{H3: Not Reached but Revealing Hierarchy}
\label{subsec:res-h3}

% [TODO — 0.5 page]
% No model reaches C-index > 0.70
% But Meditron (0.65) far above Llama (0.41-0.51) — medical specialization
% improves ordinal coherence, though the collapse limits interpretation

\subsection{H4: Supported Only on the Mistral Family}
\label{subsec:res-h4}

% [TODO — 0.5 page]
% F1@80% > F1@100% + 0.02 only for Mistral-Base A, Mistral-Base B,
% Mistral-Instruct A. Others fail (many because of collapse — abstention
% cannot help models that already always abstain).

\begin{figure}[htbp]
    \centering
    % \includegraphics[width=0.9\textwidth]{risk_coverage.pdf}
    \caption{Risk-coverage curves for the 6 Prompt-B configurations.}
    \label{fig:res-riskcov}
\end{figure}

% ----------------------------------------------------------------------------
\section{Post-hoc Recalibration and Stability}
\label{sec:res-recalibration}

% [TODO — 1.5 pages]
% Comparison table: original ECE vs temperature vs isotonic vs Platt
% Highlight isotonic: 11/12 models below 0.05 after recalibration
% 5-fold cross-validation confirms stability of gains
% Caveat: recalibration of collapsed models is trivial (calibrating a
% constant is nearly free) — not a substitute for classification quality.

\begin{finding}[F6: Post-hoc Isotonic Recalibration Works]
    Isotonic regression fitted on 70\% of the data brings 11 of 12
    configurations below the well-calibrated threshold of 0.05, robustly
    across 5-fold cross-validation. However, for collapsed models the
    ``correction'' merely calibrates a near-constant prediction and does not
    restore classification quality.
\end{finding}

% ----------------------------------------------------------------------------
\section{Conformal Prediction and Inter-Model Consensus}
\label{sec:res-conformal-consensus}

\subsection{Conformal Prediction Achieves Target Coverage}
\label{subsec:res-conformal}

% [TODO — 0.5 page]
% All models achieve 90% empirical coverage
% Set size varies: Mistral-Base B has the tightest sets (1.95 avg)

\subsection{Inter-Model Consensus}
\label{subsec:res-consensus}

% [TODO — 0.5 page]
% When all 6 models unanimous (163 cases): accuracy 75.5%
% Suggests ensembling potential (see future work)

% ----------------------------------------------------------------------------
\section{Baselines and Composite Utility Score}
\label{sec:res-baselines}

% [TODO — 1 page — NEW SECTION]
% Comparison with naïve baselines:
% - Majority ("not excluded"): 47.7% accuracy, ECE 0.00
% - Always NEI: 28.9% accuracy, ECE 0.21
%
% Key insight: several LLM configurations underperform the naïve NEI
% baseline in ECE, and their accuracy is indistinguishable from it.
%
% Composite utility score (F1_macro × (1−ECE) × diversity):
%   1. Mistral-Instruct-B (0.237)
%   2. Mistral-Instruct-A (0.173)
%   3. Mistral-Base-A (0.162)
%   4. Mistral-Base-B (0.150)
%   ...
%  12. BioMistral-B (0.000)

% ----------------------------------------------------------------------------
\section{Summary of Results}
\label{sec:res-summary}

% [TODO — 0.5 page]
% Summary table:
%   H1: supported (all 12 configurations, ECE > 0.05)
%   H2: supported (4 permutation tests, all p < 0.003 after Bonferroni)
%   H3: not reached (C-index cap at 0.65 for Meditron)
%   H4: partial (Mistral family only)
%
% 9 original findings (F1–F9)
